\subsection{Photonic composition of Bendersky--Glaisher level \texorpdfstring{\(A_2\)}{A2}}
\label{v:subsec:bendersky-fotonica}

The number density in Proposition~\ref{v:prop:densidades-termicas}
has a second exact reading that retains its spectral provenance.
Let \(\mathcal Z_3\) denote the normalized scalar functional produced by
the tritic spectral reader and subsequently recognized as the zeta function
in the Archimedean realization. Let
\[
 H_k=\sum_{j=1}^{k}\frac1j,\qquad H_0=0,
\]
and let \(\mathsf B_j\) be the Bernoulli numbers under the convention
\(\mathsf B_1=-1/2\). To avoid any collision with lattice notation,
define the Bendersky--Glaisher levels by
\begin{equation}
 A_k^{\mathrm{BG}}
 =\exp\!\left(
   \frac{H_k\mathsf B_{k+1}}{k+1}-\mathcal Z_3'(-k)
 \right),\qquad k\in\mathbb N_0.
 \label{v:eq:bendersky-bg}
\end{equation}
The Bernoulli term fixes the normalization of the regularized product; it
does not introduce a physical constant. The coordinate
\(\pi_{\mathrm{HMT}}\) entering the functional equation has already
been expressed by the enriched APP--TRIT--TPK state before this reading.

\begin{proposition}[Photonic composition at level two]
\label{v:prop:bendersky-fotonica}
In the scalar normalization \(\mathcal Z_3=\zeta\),
\begin{equation}
 \log A_2^{\mathrm{BG}}
 =-\mathcal Z_3'(-2)
 =\frac{\mathcal Z_3(3)}{4\pi_{\mathrm{HMT}}^2}.
 \label{v:eq:bendersky-nivel-dos}
\end{equation}
Consequently, in the calendar thermal chart,
\begin{align}
 n_\gamma\ell_P^3
  &=\frac{8\log A_2^{\mathrm{BG}}}{(\log3)^3}
   =-\frac{8\mathcal Z_3'(-2)}{(\log3)^3},
   \label{v:eq:numero-fotones-bendersky}\\
 \frac{u_\gamma}{n_\gamma k_B\Theta_{108}}
  &=\frac{\pi_{\mathrm{HMT}}^2}{120\log A_2^{\mathrm{BG}}},
   \label{v:eq:energia-media-bendersky}\\
 \frac{s_\gamma}{n_\gamma k_B}
  &=\frac{\pi_{\mathrm{HMT}}^2}{90\log A_2^{\mathrm{BG}}}.
   \label{v:eq:entropia-media-bendersky}
\end{align}
\end{proposition}

\begin{proof}
The completed form satisfies
\[
 \pi_{\mathrm{HMT}}^{-s/2}\Gamma(s/2)\mathcal Z_3(s)
 =\pi_{\mathrm{HMT}}^{-(1-s)/2}
  \Gamma((1-s)/2)\mathcal Z_3(1-s).
\]
At \(s=-2\), the simple zero of \(\mathcal Z_3\) cancels the pole of
\(\Gamma(s/2)\). Comparing constant terms and using
\(\Gamma(3/2)=\sqrt{\pi_{\mathrm{HMT}}}/2\) after Archimedean
recognition gives
\[
 \mathcal Z_3'(-2)=-\frac{\mathcal Z_3(3)}{4\pi_{\mathrm{HMT}}^2}.
\]
Because \(\mathsf B_3=0\), Definition~\eqref{v:eq:bendersky-bg}
gives the first equality in~\eqref{v:eq:bendersky-nivel-dos}.
Substitution of
\(\mathcal Z_3(3)=4\pi_{\mathrm{HMT}}^2\log A_2^{\mathrm{BG}}\)
into~\eqref{v:eq:densidades-calendario} proves
\eqref{v:eq:numero-fotones-bendersky}. Finally,
\(k_B\Theta_{108}=E_P/\log3\); dividing the energy and entropy
densities by the number density proves
\eqref{v:eq:energia-media-bendersky} and
\eqref{v:eq:entropia-media-bendersky}.
\end{proof}

The evaluation \(A_2^{\mathrm{BG}}=1.0309167521973921\ldots\) is a
numerical test of the three identities, not generative input. The
Bendersky--Glaisher spectral development contains the complete hierarchy of
regularized derivatives and its correspondence with odd spectral values; the present section establishes the
reciprocal composition used specifically by dimensionless photon-number
density. It neither
redefines \(k_B\), identifies a density with a regularized operator, nor
alters any of the previously proved \(\zeta(3)\) densities.
